\documentclass[11pt]{article}
\usepackage[margin=1in]{geometry}
\usepackage[hidelinks]{hyperref}
\usepackage{fancyvrb}
\title{bulk\_extractor 2.2 Current Operating Guide}
\author{bulk\_extractor project}
\date{July 2026}

\begin{document}
\maketitle

\section*{Scope}
This guide describes the supported 2.2 command-line workflow.

\section*{Build from source}
Clone the repository or download a release from
\url{https://github.com/simsong/bulk_extractor/releases}. Install the platform
dependencies described by the repository README, then run:
\begin{Verbatim}
bash bootstrap.sh
./configure
make
make check
sudo make install
\end{Verbatim}
The README is the canonical build and platform-support reference. Native
Windows builds are not supported. The repository's Windows MinGW workflow
cross-compiles a standalone \texttt{bulk\_extractor64.exe} on Ubuntu for each
pull request. Download it from the Windows workflow artifact; it is not a
signed installer or release distribution. The workflow verifies the PE imports
and runs the executable on a Windows runner against a Unicode-named input file.

\section*{Run a scan}
Choose a new output directory and an input image or file:
\begin{Verbatim}
bulk_extractor -o output mydisk.raw
\end{Verbatim}
Use \texttt{bulk\_extractor --help} on the installed binary as the authoritative
reference for options and scanner settings. A scan produces feature files,
histograms where enabled, carved content where applicable, and DFXML run
metadata in the output directory.

\section*{Develop scanners}
The maintained scanner contract is
\url{https://github.com/simsong/bulk_extractor/blob/main/doc/scanner_api.md}.
Read it before changing a built-in scanner or starting a loadable scanner, and
begin from \texttt{doc/scanner\_template.cpp}. It documents the 2.x lifecycle,
concurrency requirements, feature-recorder rules, and module factory.

\end{document}
